% ---------------------------------------------------------------- packages
\usepackage[utf8]{inputenc}
\usepackage[T1]{fontenc}
\usepackage[safe]{tipa}
\IfFileExists{lmodern.sty}{\usepackage{lmodern}\usepackage{microtype}}{\usepackage[expansion=false]{microtype}}
\usepackage[margin=2.6cm]{geometry}
\usepackage{amsmath,amssymb,amsthm,bm}
\usepackage{booktabs,tabularx,array}
\usepackage{graphicx}
\usepackage{xcolor}
\usepackage{enumitem}
\usepackage[numbers,sort&compress]{natbib}
\usepackage[colorlinks=true,linkcolor={blue!50!black},citecolor={green!35!black},urlcolor={blue!70!black},
  bookmarksnumbered=true,bookmarksopen=true,bookmarksopenlevel=1,pdfstartview=FitH]{hyperref}
\hypersetup{pdftitle={Physics of Causality: gravity as the response of a causal order},
  pdfauthor={Pavel Sikula},
  pdfsubject={Causal sets, Benincasa-Dowker action, emergent gravity, pre-registered numerical tests},
  pdfkeywords={causal set theory; Benincasa-Dowker action; emergent gravity; non-locality; post-Newtonian gamma}}
\usepackage{url}
\setlist{itemsep=2pt,topsep=4pt}
\numberwithin{equation}{section}

% ---------------------------------------------------------------- environments
\theoremstyle{definition}
\newtheorem{axiom}{Axiom}
\newtheorem{definition}{Definition}[section]
\newtheorem{rule_}{Rule}[section]
\theoremstyle{remark}
\newtheorem*{remark}{Remark}
\newtheorem*{status}{Status}

% a box for the verdict of a laboratory test
\newcommand{\verdict}[1]{\textsc{#1}}
% status label of a statement in the ontology part
\newcommand{\stat}[1]{\unskip\nobreak\hfill\mbox{\footnotesize\textsf{[#1]}}}
% editorial notes (to be removed before publication)
\newcommand{\todo}[1]{\textcolor{red}{[TODO: #1]}}

% ---------------------------------------------------------------- notation
\newcommand{\lk}{\ell_k}
\newcommand{\Pib}{\Pi}
\newcommand{\dd}{\mathrm{d}}
\newcommand{\BD}{\mathrm{BD}}
\newcommand{\Khat}{\hat K}
\newcommand{\Ex}{\mathbb{E}}
\renewcommand{\vec}[1]{\bm{#1}}
% coefficients of the static quadratic form of the action of the order (Einstein--Hilbert: 6, 8, 2)
\newcommand{\cA}{\mathcal{A}}
\newcommand{\cB}{\mathcal{B}}
\newcommand{\cC}{\mathcal{C}}
